% --------------------------------------------------
\subsubsection{Диагональ зорилтот матриц}
% --------------------------------------------------

Диагональ зорилтот матриц нь хөрөнгө тус бүрийн өгөөжийн
түүврийн вариацыг хадгалж, харин хөрөнгүүдийн өгөөж хоорондын
ковариацыг тэг гэж үзнэ. Уг зорилтот матрицыг

\[
T_{\mathrm{diag}}
=
\operatorname{diag}(S)
=
\operatorname{diag}
(s_{11},s_{22},\ldots,s_{NN})
\]

гэж тодорхойлно. Матриц хэлбэрээр

\[
T_{\mathrm{diag}}
=
\begin{pmatrix}
	s_{11} & 0      & \cdots & 0\\
	0      & s_{22} & \cdots & 0\\
	\vdots & \vdots & \ddots & \vdots\\
	0      & 0      & \cdots & s_{NN}
\end{pmatrix}.
\]

Өөрөөр хэлбэл,

\[
(T_{\mathrm{diag}})_{ij}
=
\begin{cases}
	s_{ii}, & i=j,\\
	0,      & i\neq j.
\end{cases}
\]

Иймээс диагональ зорилтот матриц нь хөрөнгө тус бүрийн
өгөөжийн түүврийн вариацыг өөрчлөхгүй хадгалах боловч
хөрөнгүүдийн өгөөж хоорондын ковариацыг тэг болгон
хялбаршуулдаг.

Диагональ зорилтот матрицад суурилсан шугаман агшилтын
үнэлэгчийг

\[
\widehat{\Sigma}_{\mathrm{diag}}
=
(1-\widehat{\alpha}_{\mathrm{diag}})S
+
\widehat{\alpha}_{\mathrm{diag}}T_{\mathrm{diag}},
\qquad
0\leq\widehat{\alpha}_{\mathrm{diag}}\leq1,
\]

гэж тодорхойлно. Энд
$\widehat{\alpha}_{\mathrm{diag}}$ нь диагональ зорилтот
матрицад харгалзах агшилтын эрчим юм.



\paragraph{Агшилтын эрчмийн үнэлгээ.}

Диагональ зорилтот матриц нь түүврийн ковариацын
матрицын диагональ элементүүдийг өөрчлөхгүй
хадгалдаг тул агшилт нь зөвхөн диагоналийн бус
элементүүдэд үйлчилнэ.

Иймээс агшилтын эрчмийг үнэлэхдээ
диагоналийн бус элементүүдийн квадрат
алдааг авч үзнэ.

Түүврийн ковариацын матриц болон
диагональ зорилтот матрицын хоорондын
нормчилсон квадрат Фробениусын зайг
\[
\widehat d_{\mathrm{diag}}
=
\|S-T_{\mathrm{diag}}\|^2
=
\frac{1}{N}
\sum_{\substack{i,j=1\\i\neq j}}^{N}
s_{ij}^{2}
\]
гэж тодорхойлно.

Учир нь $S$ болон $T_{\mathrm{diag}}$
матрицын диагональ элементүүд ижил тул
тэдгээрийн ялгавар диагональ дээр тэг байна.

Өмнөх тогтмол вариацын зорилтот матрицын хэсэгт үзсэнтэй адилтаар $\widehat{b}_{diag}$-г дараах байдлаар тодорхойлно. 
\[
\widehat b_{\mathrm{diag}}
	=
	\frac{1}{nN}
	\sum_{\substack{i,j=1\\i\neq j}}^{N}
	\left[
	\frac{1}{n}
	\sum_{t=1}^{n}
	(z_{tij}-s_{ij})^2
	\right].
\]
Иймээс диагональ зорилтот матрицад
харгалзах агшилтын эрчмийн
практик үнэлгээг
\[
	\widehat{\alpha}_{\mathrm{diag}}
	=
	\max\left\{
	0,\,
	\min\left\{
	1,\,
	\frac{\widehat b_{\mathrm{diag}}}
	{\widehat d_{\mathrm{diag}}}
	\right\}
	\right\}
\]
гэж байгуулна.

Эцэст нь агшилтын ковариацын
матрицыг
\[
\widehat{\Sigma}_{\mathrm{diag}}
=
(1-\widehat{\alpha}_{\mathrm{diag}})S
+
\widehat{\alpha}_{\mathrm{diag}}
T_{\mathrm{diag}}
\]
гэж тооцно.

\paragraph{Тооцоолох алгоритм.}

Диагональ зорилтот матрицад суурилсан
агшилтын ковариацын матрицыг дараах
дарааллаар тооцоолно.

\begin{enumerate}
\item Диагональ зорилтот матрицыг байгуулах.

Өмнөх хэсэгт тооцсон түүврийн ковариацын
матриц $S=(s_{ij})$-ийг ашиглан
\[
T_{\mathrm{diag}}
=
\operatorname{diag}
(s_{11},s_{22},\ldots,s_{NN})
\]
гэж байгуулна.
	\item Диагоналийн бус элементүүдийн
	квадрат алдааг үнэлэх.
	\[
	\widehat b_{\mathrm{diag}}
	=
	\frac{1}{nN}
	\sum_{\substack{i,j=1\\i\neq j}}^{N}
	\left[
	\frac{1}{n}
	\sum_{t=1}^{n}
	(z_{tij}-s_{ij})^2
	\right].
	\]
	\item Зорилтот матрицаас хазайх
	квадрат зайг тооцох.
	
	\[
	\widehat d_{\mathrm{diag}}
	=
	\frac{1}{N}
	\sum_{\substack{i,j=1\\i\neq j}}^{N}
	s_{ij}^{2}.
	\]
	
	\item Агшилтын эрчмийг үнэлэх.
	
	$\widehat d_{\mathrm{diag}}>0$ үед
	\[
	\widehat{\alpha}_{\mathrm{diag}}
	=
	\max\left\{
	0,\,
	\min\left\{
	1,\,
	\frac{\widehat b_{\mathrm{diag}}}
	{\widehat d_{\mathrm{diag}}}
	\right\}
	\right\}.
	\]
	
	\item Агшилтын ковариацын
	матрицыг байгуулах.
	
	\[
		\widehat{\Sigma}_{\mathrm{diag}}
		=
		(1-\widehat{\alpha}_{\mathrm{diag}})S
		+
		\widehat{\alpha}_{\mathrm{diag}}
		T_{\mathrm{diag}}.
	\]
\end{enumerate}
